\chapter{Introduction}
\label{chap:introduction}

\section{Background and Motivation}
\label{sec:intro_motivation}

Computed tomography (\gls{ct}) is one of the most widely used imaging modalities in modern medicine. By acquiring X-ray attenuation measurements from numerous projection angles and reconstructing them into cross-sectional images, CT enables non-invasive visualization of internal anatomy with high diagnostic accuracy. This capability supports the detection of pulmonary nodules, the characterization of cardiovascular lesions, the staging of oncologic disease, and the evaluation of trauma without the need for surgical intervention \cite{Sadia2024, Komolafe2024}. Due to its substantial clinical utility, the global volume of CT examinations has increased steadily over the past two decades.

This growth, however, is accompanied by an unavoidable cost: exposure to ionizing radiation. Each CT scan delivers an X-ray dose to the patient, and cumulative exposure is associated with a small but well-documented risk of stochastic biological harm, including radiation-induced malignancy \cite{Sadia2024, AAPM2016}. In particular, a large retrospective cohort study reported an increased risk of leukemia and brain tumors in children who underwent repeated CT examinations~\cite{Pearce2012}. Because the radiation dose of the patient scales directly with the current of the X-ray tube, the most straightforward strategy for reducing the dose is to lower the current of the tube, which forms the basis of \gls{ldct}. However, this approach introduces a substantial trade-off. As fewer photons reach the detector, the measured sinogram becomes increasingly dominated by quantum noise governed by Poisson statistics. When reconstructed using \gls{fbp}, which remains the standard analytical method in clinical practice, this noise propagates into the image domain as low-frequency granular noise, streak artifacts, and reduced soft-tissue contrast~\cite{Leuschner2021}. As a result, image quality can be diagnostically compromised, thus limiting the extent to which the radiation dose can be reduced in routine practice~\cite{AAPM2016}.

Noise is not the only source of image degradation in CT. In routine clinical imaging, CT scans are often affected by multiple co-occurring physical artifacts that arise from the acquisition process itself rather than from photon statistics alone~\cite{Boas2012}. Even minor patient motion during scanning, such as breathing, swallowing, or involuntary cardiac movement, can introduce inconsistencies across projections, leading to blurring and ghosting in the reconstructed image, particularly around moving structures \cite{LuMackie2002, Lu2005}. Detector miscalibration or damage to individual detector elements may further produce ring artifacts, which appear as concentric circular bands centered on the axis of rotation and can obscure subtle lesions near the center of the field of view \cite{Sijbers2004, Prell2009}. In addition, metallic implants, including hip prostheses, dental fillings, pacemaker leads, and surgical clips, cause extreme attenuation that leads to photon starvation along affected projection paths and generate pronounced bright and dark streak artifacts, often rendering substantial image regions uninterpretable \cite{ZhangYu2018, Gjesteby2016}. Importantly, these degradations rarely occur in isolation. For example, a low dose postoperative patient may simultaneously present with Poisson noise, ring artifacts caused by detector drift, and metal artifacts associated with the implant. Nevertheless, most deep learning studies continue to address these sources of degradation independently rather than as interacting components of a single clinical imaging problem \cite{Liang2025}.

During the past three decades, substantial algorithmic effort has been devoted to recovering diagnostic image quality from degraded CT data. Classical approaches include hand-crafted image-domain post-processing filters, total variation minimization, dictionary-learning-based denoising, and non-local means methods \cite{Sadia2024, Tian2020}. Among these, Block-Matching and 3D Filtering (\gls{bm3d}) \cite{Dabov2007} occupies a distinctive position because it exploits non-local self-similarity to achieve state-of-the-art performance for additive Gaussian noise without requiring training data, making it a natural benchmark against which learned methods are often evaluated. On the reconstruction side, iterative methods that incorporate statistical noise models and anatomical regularization priors, most notably the constrained total-variation minimization framework of Sidky and Pan \cite{SidkyPan2008}, can outperform FBP in low-dose settings, although their computational cost remains a significant obstacle to routine clinical use \cite{Beister2012}.

Supervised deep learning has emerged as the dominant paradigm in CT image restoration. Convolutional encoder--decoder architectures, popularized in biomedical imaging by the U-Net~\cite{Ronneberger2015}, have been widely adapted for CT restoration tasks. Trained on paired degraded and reference images, these models learn an end-to-end mapping from corrupted input to high-quality output. A landmark example is the Residual Encoder--Decoder CNN (RED-CNN) \cite{Chen2017}, which introduced skip connections between the corresponding encoder and decoder layers to preserve fine structural detail during compression. In parallel, DnCNN \cite{Zhang2017} showed that learning residual noise, rather than directly predicting clean image, can accelerate convergence and improve denoising performance, following a strategy inspired by residual learning in ResNet-style architectures \cite{He2016}. A related line of work further demonstrated that the same encoder--decoder principle can be applied to the output of an analytic reconstructor in general inverse problems \cite{Jin2017}, while adversarial and perceptual losses were later introduced to mitigate the over smoothing often associated with pure $\ell^{2}$ optimization \cite{Yang2018WGAN}. These ideas have since been combined and extended in numerous forms, including attention-based architectures \cite{Komolafe2024, Zhang2023}, dual-domain methods that jointly process sinogram and image data~\cite{Wang2020, Zhang2022, Lin2019DuDoNet}, and model-based unrolled networks that explicitly incorporate the forward operator into the reconstruction architecture~\cite{Aggarwal2019MoDL, Hammernik2018VarNet}. Recent benchmarking studies have also shown that the performance reported of these methods can be highly sensitive to the corruption protocol and evaluation data set, underscoring the need for reproducible and physics-consistent test conditions~\cite{Kiss2024}.

Despite this rapid progress, two important questions remain insufficiently explored. First, most existing methods are designed and evaluated for a single isolated degradation type, most commonly simulated Poisson noise at a prescribed dose level, while implicitly assuming the absence of other artifact sources. Although this simplification facilitates benchmarking, it creates a clear mismatch with clinical reality, where multiple degradation mechanisms often coexist within the same scan~\cite{Boas2012, Liang2025}. At present, there is no widely established protocol for jointly simulating this compound corruption distribution in the sinogram domain in a physics-consistent manner. Second, although cascade architectures, in which two or more restoration networks are applied sequentially, are widely used in practice, the question of how the stages should be coupled has received comparatively little systematic study. In particular, it remains unclear whether stages should be trained independently with gradients stopped at stage boundaries, or instead optimized jointly in an end-to-end manner. Similarly, the effect of formulating each stage as a residual corrector rather than a complete restoration module remains poorly understood. These design choices have important implications for gradient propagation, optimization stability, and the extent to which performance can be improved by increasing cascade depth.

This thesis addresses both questions directly by using the LoDoPaB-CT benchmark \cite{Leuschner2021} as its experimental testbed. LoDoPaB-CT is a large-scale, open-access dataset of simulated parallel-beam CT measurements derived from the LIDC/IDRI lung CT database.

\section{Problem Statement}
\label{sec:intro_problem}

This thesis addresses a twofold central challenge. First, it asks how realistic multi-artifact degradation can be simulated faithfully in the sinogram domain so that learned models are exposed to compound corruption distributions that better reflect clinical acquisition. Second, under such a corruption regime, it investigates how architectural depth, specifically the choice between independent-stage, end-to-end, and residual cascade formulations of encoder--decoder networks, influences restoration quality and whether a meaningful performance ceiling exists for additional image-domain stages. These two questions are inherently linked. The effect of cascade depth becomes substantively interesting only when the corruption is sufficiently complex that a single-stage model may be inadequate. Conversely, the value of a realistic corruption protocol can only be assessed through a well-controlled comparison of architectural alternatives in which depth is isolated as the independent variable. Together, these questions define a systematic study of image-domain restoration capacity under physics-informed corruption.

\section{Contributions}
\label{sec:intro_contributions}

This thesis makes two primary contributions:

\begin{enumerate}
    \item \textbf{A physics-informed multi-artifact sinogram corruption protocol.} This thesis develops and validates a modular, reproducible pipeline for the joint simulation of Poisson quantum noise, motion artifacts, ring artifacts, and metal artifacts in the sinogram domain, each at controllable levels of severity. Each component is designed to reflect the physical mechanism that gives rise to the corresponding degradation in real CT acquisition: Poisson photon statistics for quantum noise~\cite{Leuschner2021}, projection inconsistencies induced by motion~\cite{LuMackie2002, Lu2005}, per-detector gain perturbations for ring artifacts~\cite{Sijbers2004, Prell2009}, and high-attenuation insertions with associated photon starvation for metal artifacts~\cite{ZhangYu2018, Gjesteby2016}. The protocol is parameterized to match the parallel-beam geometry and photon statistics of LoDoPaB-CT~\cite{Leuschner2021}, and individual artifact types are sampled independently to generate a compound corruption distribution spanning a broad and clinically plausible range of degradation conditions. Crucially, all corruptions are introduced prior to \gls{fbp} reconstruction, thereby preserving physical consistency in the way they propagate into the image domain.

    \item \textbf{A systematic study of cascade architectures.} Three cascade formulations of a U-Net encoder--decoder are trained and evaluated under the multi-artifact corruption protocol described above: (i)~an \emph{independent-stage} cascade, in which both stages are trained simultaneously but a stop-gradient operator at the stage boundary prevents the Stage-2 loss from reaching Stage-1 parameters, so that each stage is updated only by its own loss; (ii)~an \emph{end-to-end cascade}, in which all stages are optimized jointly under deep supervision~\cite{Lee2015DSN}, with an auxiliary loss applied to each intermediate stage output in addition to the final reconstruction loss, thereby allowing gradients to propagate through the full cascade while providing explicit supervision at every stage; and (iii)~a \emph{residual cascade}, in which each stage predicts an additive correction to the output of the previous stage rather than a full restored image. These three formulations are compared against three baselines --- \gls{bm3d}~\cite{Dabov2007} applied to the FBP reconstruction, a single-stage U-Net~\cite{Ronneberger2015}, and \gls{redcnn}~\cite{Chen2017} --- using \gls{psnr} and \gls{ssim} as the primary evaluation metrics, consistent with current CT benchmarking practice~\cite{Kiss2024}.
\end{enumerate}

The central finding of this thesis is that a single well-designed encoder--decoder network already operates close to the performance ceiling attainable in the image domain under the compound corruption conditions considered. Adding further image-domain cascade stages yields only limited benefit: the naive independent-stage cascade improves on the matched single-stage baseline by just \(+0.15\)\,dB PSNR, while the end-to-end cascade, despite joint training with gradient flow across stages, achieves only \(+0.46\)\,dB and therefore only marginally surpasses the naive variant. These diminishing returns are supported by gradient-norm analysis, which indicates that Stage~2 receives a substantially weaker learning signal once Stage~1 has converged. The residual cascade recovers a small but consistent gain of \(+0.54\)\,dB relative to the naive independent-stage baseline, and, when combined with dual-filter conditioning, is the only cascade configuration that matches the larger single-stage U-Net. Rather than constituting a negative result, this architectural ceiling is scientifically informative: it helps to localize the image-domain information bottleneck and thereby motivates dual-domain approaches that introduce additional information from the sinogram as the logical next step~\cite{Wang2020, Zhang2022}.

\section{Thesis Outline}
\label{sec:intro_outline}

The remainder of this thesis is organized as follows. \Cref{chap:background} provides the theoretical background on CT imaging, the mechanisms of low-dose degradation and physical artifacts, and the deep learning literature relevant to CT image restoration. \Cref{chap:methodology} describes the methodological framework of the thesis, including the physics-informed multi-artifact sinogram corruption protocol (\Cref{chap:corruption}) and the cascaded encoder--decoder architectures, comprising the single-stage baseline, the alternative-filter input scheme, the three coupling strategies, and the common training procedure (\Cref{chap:architectures}). \Cref{chap:experiments} presents the quantitative and qualitative results, together with the gradient analysis and a critical discussion of the findings. Finally, \Cref{chap:conclusion} summarizes the main contributions, discusses the limitations of the study, and outlines directions for future research.
